% Ch16 self-study -- "I do / You do" ladder for Zumdahl 16.1 plus the
% examinable slice of 16.2. Nine ladders + a ten-question mixed set =
% 46 questions, eight figures.
%
% SCOPE, re-verified against the CED PDF on 2026-08-23 (see
% TEXTBOOK-MAP): 16.1 is fully examinable (CED 7.11, 7.12, 8.11).
% 16.2 is PARTIALLY examinable -- predicting whether a precipitate
% forms by comparing Q with Ksp is in (the CED's own 7.11 activity has
% students do exactly this); selective precipitation and the
% qualitative-analysis schemes are out ("qualitative analysis" appears
% nowhere in the CED). 16.3 is out ("complex ion" appears nowhere).
% The 8.11 exclusion is honoured: pH-and-solubility stays qualitative.
\documentclass{apchem-worksheet}

\apcunitword{Chapter}
\apcunit{16}
\apcunittitle{Solubility Equilibria}
\apcdoctitle{Self-Study \textbullet\ Chapter 16, I~do / You~do}
\apcsubtitle{Zumdahl \S16.1 $+$ \S16.2 (Q vs $K_{sp}$) \textbullet\ PDF pp.\ 805--819}

\begin{document}
\apctitleblock

\begin{apcnote}[How to use these notes --- read this first]
Each skill is a \textbf{ladder}: a fully worked example in the
solid-framed box, then \textbf{four} YOUR TURN questions in the dashed
box --- same skill, new numbers, no help. Work all four before checking
the gray \emph{check:} line; a tracker tick needs all four right first
time. A ten-question \textbf{mixed set} follows Ladder 9.

\textbf{What is in and what is out --- checked against the CED
itself.} \S16.1 ($K_{sp}$, molar solubility, the common-ion effect, pH
and solubility) is fully assessed: CED topics 7.11, 7.12 and 8.11.
From \S16.2, predicting whether a \textbf{precipitate forms} by
comparing $Q$ with $K_{sp}$ is assessed --- it is Chapter 13's $Q$
machinery applied to dissolving. The rest of \S16.2 (selective
precipitation schemes, qualitative analysis) and all of \S16.3
(complex ions) do not appear anywhere in the CED and are not treated
here.

\textbf{What you must bring with you.} Chapter 13, again. A $K_{sp}$
problem is an equilibrium problem where the ICE table is so simple it
barely deserves the name --- the initial concentrations are zero and
everything is written in one unknown. If heterogeneous equilibria
(Chapter 13, Ladder 4) and $Q$ versus $K$ (Ladder 6) are solid, this
chapter is two ideas and some algebra.

\textbf{The one idea underneath everything.} A ``saturated'' solution
is not a solution where dissolving has stopped --- it is a
\textbf{dynamic equilibrium} where dissolving and precipitating run at
equal rates, and $K_{sp}$ is just the equilibrium constant for it.
\end{apcnote}

% =====================================================================
\section*{\sffamily Ladder 1 \textbullet\ The dissolution equilibrium
and $K_{sp}$}
\zsec{16.1}

Drop solid \ce{AgCl} into water. A little dissolves, then the solution
saturates --- and from that moment ions leave the crystal exactly as
fast as they return to it.

\begin{center}
\begin{tikzpicture}[font=\sffamily\footnotesize]
  % saturated solution beaker: fill BEFORE stroke
  \fill[apcgray!22] (0.6,0.4) rectangle (3.6,2.6);
  \draw[very thick] (0.6,3.0) -- (0.6,0.4) -- (3.6,0.4) -- (3.6,3.0);
  % solid at the bottom
  \fill[apcgray!75] (1.0,0.4) rectangle (1.9,0.72);
  \fill[apcgray!75] (2.2,0.4) rectangle (3.1,0.65);
  % dissolved ions
  \foreach \x/\y/\l in {1.0/1.6/{\ce{Ag+}}, 2.0/2.1/{\ce{Cl-}},
                        2.9/1.5/{\ce{Ag+}}, 1.5/2.35/{\ce{Cl-}},
                        2.5/1.05/{\ce{Cl-}}, 1.6/1.1/{\ce{Ag+}}} {
    \node[font=\tiny] at (\x,\y) {\l};}
  % exchange arrows
  \draw[-{Stealth[length=2mm]},thick] (1.35,0.78) -- (1.35,1.25);
  \draw[{Stealth[length=2mm]}-,thick] (2.75,0.72) -- (2.75,1.2);
  \node[align=center] at (2.1,-0.25)
    {dissolving and precipitating\\at \textbf{equal rates}};

  \node[anchor=west,align=left] at (4.6,2.3)
    {$\ce{AgCl(s) <=> Ag+(aq) + Cl-(aq)}$};
  \node[anchor=west,align=left] at (4.6,1.5)
    {$K_{sp} = [\ce{Ag+}][\ce{Cl-}]$};
  \node[anchor=west,align=left,font=\sffamily\footnotesize] at (4.6,0.7)
    {the solid is \textbf{omitted} --- Chapter 13, Ladder 4:\\
     a pure solid never appears in an equilibrium expression};
\end{tikzpicture}
\end{center}

\begin{apctrap}
$K_{sp}$ expressions have \textbf{no denominator}. The only thing on
the left of the dissolution equation is a pure solid, and pure solids
are omitted --- so nothing is left to divide by. Writing
$[\ce{Ag+}][\ce{Cl-}]/[\ce{AgCl}]$ is the error this ladder exists to
prevent.
\end{apctrap}

\begin{workedexample}[I do: three expressions, mind the exponents]
\textbf{(a) $\ce{BaSO4(s) <=> Ba^2+(aq) + SO4^2-(aq)}$}
\[ K_{sp} = [\ce{Ba^2+}][\ce{SO4^2-}] \]

\textbf{(b) $\ce{PbI2(s) <=> Pb^2+(aq) + 2I-(aq)}$}
\[ K_{sp} = [\ce{Pb^2+}][\ce{I-}]^{2} \]
The coefficient 2 on iodide becomes an exponent --- exactly as in every
equilibrium expression since Chapter 13.

\textbf{(c) $\ce{Ca3(PO4)2(s) <=> 3Ca^2+(aq) + 2PO4^3-(aq)}$}
\[ K_{sp} = [\ce{Ca^2+}]^{3}[\ce{PO4^3-}]^{2} \]

\textbf{The procedure is always the same two steps.} Write the
dissolution equation --- solid on the left, ions on the right, with the
subscripts of the formula becoming the coefficients of the ions. Then
write the product of the ion concentrations, each raised to its
coefficient. There is nothing else to it, and there is no denominator.

\textbf{Why $K_{sp}$ values are so small.} These are all salts your
solubility rules call ``insoluble''. Insoluble does not mean zero
dissolves --- it means the equilibrium lies far to the left, which is
exactly the statement $K_{sp} \ll 1$. The solubility rules of Chapter 4
are the qualitative shadow of this chapter's numbers: $K_{sp} > 1$
corresponds to the salts the rules call soluble.
\end{workedexample}

\begin{yourturn}[YOUR TURN 1 --- four questions]
Write the dissolution equation and the $K_{sp}$ expression.
\begin{enumerate}[leftmargin=*,itemsep=0.5em]
  \item \ce{AgBr} \workspace[1.6cm]
  \item \ce{CaF2} \workspace[1.6cm]
  \item \ce{Ag2CrO4} \workspace[1.8cm]
  \item Why does the solid not appear in the expression?
        \workspace[1.6cm]
\end{enumerate}
\selfcheck{(a) $[\ce{Ag+}][\ce{Br-}]$ \quad (b)
$[\ce{Ca^2+}][\ce{F-}]^{2}$ \quad (c)
$[\ce{Ag+}]^{2}[\ce{CrO4^2-}]$ \quad (d) pure solids are omitted}
\keyonly{\begin{apcanswer}(a) $\ce{AgBr(s) <=> Ag+(aq) + Br-(aq)}$;
$K_{sp} = [\ce{Ag+}][\ce{Br-}]$. (b)
$\ce{CaF2(s) <=> Ca^2+(aq) + 2F-(aq)}$;
$K_{sp} = [\ce{Ca^2+}][\ce{F-}]^{2}$. (c)
$\ce{Ag2CrO4(s) <=> 2Ag+(aq) + CrO4^2-(aq)}$;
$K_{sp} = [\ce{Ag+}]^{2}[\ce{CrO4^2-}]$. (d) Because a \textbf{pure
solid's concentration is fixed} (density over molar mass) and cannot
change --- constants are folded into $K$, so the solid is omitted. Same
rule, same reason, as Chapter 13 Ladder 4.\end{apcanswer}}
\end{yourturn}

% =====================================================================
\section*{\sffamily Ladder 2 \textbullet\ Molar solubility --- one
unknown does everything}
\zsec{16.1}

\textbf{Molar solubility}, $s$, is the number of moles of the solid
that dissolve per litre of saturated solution. Every ion concentration
is a multiple of it.

\begin{center}
\begin{tikzpicture}[font=\sffamily\footnotesize]
  \node[font=\sffamily\small\bfseries] at (5.5,3.1)
    {let $s$ mol/L dissolve --- the stoichiometry does the rest};
  \foreach \y/\form/\ions/\expr in {%
      2.35/{\ce{AgCl}}/{$[\ce{Ag+}] = s$, \ $[\ce{Cl-}] = s$}/{$K_{sp}
        = s^{2}$},
      1.6/{\ce{PbI2}}/{$[\ce{Pb^2+}] = s$, \ $[\ce{I-}] =
        \mathbf{2s}$}/{$K_{sp} = s(2s)^{2} = 4s^{3}$},
      0.85/{\ce{Ag2CrO4}}/{$[\ce{Ag+}] = \mathbf{2s}$, \
        $[\ce{CrO4^2-}] = s$}/{$K_{sp} = (2s)^{2}s = 4s^{3}$}} {
    \node[anchor=west,font=\sffamily\bfseries] at (0.4,\y) {\form};
    \node[anchor=west] at (2.1,\y) {\ions};
    \node[anchor=west] at (7.3,\y) {\expr};}
  \draw[apcgray] (0.3,2.75) -- (10.8,2.75);
  \draw[apcgray] (0.3,0.5) -- (10.8,0.5);
  \node[anchor=west] at (0.4,0.05)
    {the 2 appears \textbf{twice} for \ce{PbI2}: once as the
     coefficient ($2s$), once as the exponent ($^{2}$)};
\end{tikzpicture}
\end{center}

\begin{apctrap}
\textbf{``Doubled AND squared'' feels like counting twice. It is not.}
$[\ce{I-}] = 2s$ because each formula unit releases two iodides ---
that is stoichiometry. $[\ce{I-}]$ is then \emph{squared} because the
equilibrium expression demands the coefficient as an exponent --- that
is the law of mass action. Two different rules, applied once each.
Dropping either one is the most common $K_{sp}$ error on the exam.
\end{apctrap}

\begin{workedexample}[I do: from measured solubility to $K_{sp}$]
\textbf{(a) The molar solubility of \ce{AgCl} is
\SI{1.3e-5}{\Molar} at \SI{25}{\celsius}. Find $K_{sp}$.}

Each dissolved formula unit gives one \ce{Ag+} and one \ce{Cl-}:
\[ [\ce{Ag+}] = [\ce{Cl-}] = s = \num{1.3e-5} \]
\[ K_{sp} = s^{2} = (\num{1.3e-5})^{2} = \mathbf{\num{1.7e-10}} \]

\textbf{(b) The molar solubility of \ce{PbI2} is
\SI{1.5e-3}{\Molar}. Find $K_{sp}$.}

Here $[\ce{Pb^2+}] = s$ and $[\ce{I-}] = 2s = \num{3.0e-3}$:
\[ K_{sp} = s(2s)^{2} = 4s^{3}
          = 4(\num{1.5e-3})^{3}
          = 4 \times \num{3.375e-9}
          = \mathbf{\num{1.4e-8}} \]

\textbf{Check the setup before the arithmetic.} The commonest wrong
answer here is $s \times (2s)$ without the square, or
$s \times s^{2}$ without the doubling. Write the two ion
concentrations explicitly first --- $s$ and $2s$ --- then substitute
into the $K_{sp}$ expression you wrote in Ladder 1. Setup first,
numbers second.

\textbf{Where these measurements come from.} Evaporate a litre of the
saturated solution and weigh what is left, or measure the conductivity,
or use a spectrometer --- molar solubility is an experimental quantity,
and $K_{sp}$ tables are built from exactly this calculation.
\end{workedexample}

\begin{yourturn}[YOUR TURN 2 --- four questions]
\begin{enumerate}[leftmargin=*,itemsep=0.5em]
  \item The molar solubility of \ce{AgBr} is \SI{7.1e-7}{\Molar}.
        Find $K_{sp}$. \workspace[1.8cm]
  \item The molar solubility of \ce{PbF2} is \SI{2.0e-3}{\Molar}.
        Find $K_{sp}$. \workspace[2.0cm]
  \item For \ce{CaF2} dissolving with molar solubility $s$, give
        $[\ce{Ca^2+}]$ and $[\ce{F-}]$ in terms of $s$.
        \workspace[1.6cm]
  \item In the \ce{PbF2} calculation, why is the fluoride concentration
        both doubled and squared? \workspace[1.8cm]
\end{enumerate}
\selfcheck{(a) \num{5.0e-13} \quad (b) \num{3.2e-8} \quad (c) $s$ and
$2s$ \quad (d) stoichiometry doubles it; the mass-action law squares it}
\keyonly{\begin{apcanswer}(a) $K_{sp} = s^{2} = (\num{7.1e-7})^{2} =
\mathbf{\num{5.0e-13}}$. (b) $K_{sp} = 4s^{3} = 4(\num{2.0e-3})^{3} =
4 \times \num{8.0e-9} = \mathbf{\num{3.2e-8}}$. (c) $[\ce{Ca^2+}] =
\mathbf{s}$ and $[\ce{F-}] = \mathbf{2s}$. (d) Two separate rules:
\textbf{stoichiometry} says each formula unit releases two fluorides,
so $[\ce{F-}] = 2s$; the \textbf{equilibrium expression} raises
$[\ce{F-}]$ to the power of its coefficient, so that $2s$ is then
squared. Applying each rule once gives $4s^{3}$; skipping either gives
a wrong answer.\end{apcanswer}}
\end{yourturn}

% =====================================================================
\section*{\sffamily Ladder 3 \textbullet\ From $K_{sp}$ to solubility}
\zsec{16.1}

The reverse direction: given the tabulated constant, how much actually
dissolves?

\begin{workedexample}[I do: two stoichiometries, two algebra shapes]
\textbf{(a) $K_{sp}(\ce{BaSO4}) = \num{1.5e-9}$. Find the molar
solubility.}

1:1 salt, so $K_{sp} = s^{2}$:
\[ s = \sqrt{\num{1.5e-9}} = \mathbf{\SI{3.9e-5}{\Molar}} \]

\textbf{(b) $K_{sp}(\ce{CaF2}) = \num{4.0e-11}$. Find the molar
solubility.}

1:2 salt, so $K_{sp} = 4s^{3}$:
\[ s^{3} = \frac{\num{4.0e-11}}{4} = \num{1.0e-11}
   \qquad
   s = \sqrt[3]{\num{1.0e-11}} = \mathbf{\SI{2.2e-4}{\Molar}} \]

\textbf{The no-calculator cube root.} $\num{1.0e-11} =
\num{10e-12}$, and $\sqrt[3]{\num{10e-12}} = \sqrt[3]{10} \times
10^{-4} \approx 2.2 \times 10^{-4}$. Rewriting so the exponent is a
multiple of 3 is the whole trick, and MCQ items are built to allow it.

\textbf{Convert to grams if asked.} For \ce{BaSO4}
($M = \SI{233.4}{\gram\per\mole}$):
\[ \num{3.9e-5} \times 233.4 \approx \SI{9.0e-3}{\gram\per\litre} \]
--- about nine milligrams dissolving in a litre, which is why barium
sulfate can be swallowed for X-ray imaging even though barium ion is
toxic: the $K_{sp}$ keeps almost all of it solid.

\textbf{And note which direction is which.} Ladder 2 went measurement
$\to$ constant. This ladder goes constant $\to$ prediction. The exam
asks both, in both stoichiometries --- four variants of one skill.
\end{workedexample}

\begin{yourturn}[YOUR TURN 3 --- four questions]
\begin{enumerate}[leftmargin=*,itemsep=0.5em]
  \item A 1:1 salt has $K_{sp} = \num{4.0e-12}$. Find its molar
        solubility. \workspace[1.8cm]
  \item A 1:2 salt \ce{MX2} has $K_{sp} = \num{3.2e-11}$. Find its
        molar solubility. \workspace[2.0cm]
  \item For the salt in (a), find $[\ce{M+}]$ in the saturated
        solution. \workspace[1.4cm]
  \item Which algebra shape --- $s^{2}$ or $4s^{3}$ --- goes with which
        formula type, and what decides it? \workspace[1.8cm]
\end{enumerate}
\selfcheck{(a) \SI{2.0e-6}{\Molar} \quad (b) \SI{2.0e-4}{\Molar} \quad
(c) \SI{2.0e-6}{\Molar} \quad (d) 1:1 gives $s^{2}$; 1:2 gives
$4s^{3}$ --- the formula's stoichiometry}
\keyonly{\begin{apcanswer}(a) $s = \sqrt{\num{4.0e-12}} =
\mathbf{\SI{2.0e-6}{\Molar}}$. (b) $s^{3} = \num{3.2e-11}/4 =
\num{8.0e-12}$, so $s = \sqrt[3]{\num{8.0e-12}} =
\mathbf{\SI{2.0e-4}{\Molar}}$ --- note $\num{8.0e-12}$ was chosen to
have an exact cube root. (c) A 1:1 salt gives one cation per formula
unit, so $[\ce{M+}] = s = \mathbf{\SI{2.0e-6}{\Molar}}$. (d) The
\textbf{formula's stoichiometry} decides: a 1:1 salt (\ce{AgCl},
\ce{BaSO4}) gives $K_{sp} = s^{2}$; a 1:2 or 2:1 salt (\ce{CaF2},
\ce{Ag2CrO4}) gives $K_{sp} = 4s^{3}$. Write the ion concentrations in
terms of $s$ first and the shape falls out.\end{apcanswer}}
\end{yourturn}

% =====================================================================
\section*{\sffamily Ladder 4 \textbullet\ Comparing solubilities ---
the trap in the tables}
\zsec{16.1}

``Which salt is more soluble?'' sounds like ``which $K_{sp}$ is
bigger?'' --- and \emph{that is only true when the formulas have the
same shape.}

\begin{center}
\begin{tikzpicture}[font=\sffamily\footnotesize]
  \node[font=\sffamily\small\bfseries] at (5.4,3.3)
    {smaller $K_{sp}$, yet MORE soluble};
  % AgCl bar
  \node[anchor=west,font=\sffamily\bfseries] at (0.5,2.6) {\ce{AgCl}};
  \node[anchor=west] at (2.3,2.6) {$K_{sp} = \num{1.6e-10}$};
  \node[anchor=west] at (5.6,2.6) {$s = \num{1.3e-5}$};
  \fill[apcgray!55] (7.9,2.45) rectangle (8.36,2.75);
  % Ag2CrO4 bar
  \node[anchor=west,font=\sffamily\bfseries] at (0.5,1.85)
    {\ce{Ag2CrO4}};
  \node[anchor=west] at (2.3,1.85) {$K_{sp} = \num{9.0e-12}$};
  \node[anchor=west] at (5.6,1.85) {$s = \num{1.3e-4}$};
  \fill[apcgray!55] (7.9,1.7) rectangle (10.5,2.0);
  \node[font=\sffamily\footnotesize,anchor=west] at (7.9,1.35)
    {molar solubility, to scale ($\times 10$)};
  \draw[apcgray] (0.4,1.0) -- (10.8,1.0);
  \node[anchor=west] at (0.5,0.55)
    {\ce{Ag2CrO4}'s $K_{sp}$ is \textbf{18 times smaller}, but its
     molar solubility is \textbf{10 times larger} ---};
  \node[anchor=west] at (0.5,0.05)
    {because $s^{2}$ and $4s^{3}$ are different functions, the two
     constants live on different scales};
\end{tikzpicture}
\end{center}

\begin{workedexample}[I do: when you may compare, and when you must
calculate]
\textbf{Rule 1 --- same formula type: compare $K_{sp}$ directly.}
\ce{AgCl} ($\num{1.6e-10}$), \ce{AgBr} ($\num{5.0e-13}$), \ce{AgI}
($\num{1.5e-16}$) are all 1:1, so the ranking of $K_{sp}$ \emph{is} the
ranking of solubility: $\ce{AgCl} > \ce{AgBr} > \ce{AgI}$. No
calculation needed --- and this is CED 7.11.A.2's ``predict the
relative solubility'' skill.

\textbf{Rule 2 --- different formula types: you must compute $s$ for
each.} Compare \ce{AgCl} with \ce{Ag2CrO4}:
\[ s(\ce{AgCl}) = \sqrt{\num{1.6e-10}} = \num{1.3e-5} \]
\[ s(\ce{Ag2CrO4}) = \sqrt[3]{\frac{\num{9.0e-12}}{4}}
   = \sqrt[3]{\num{2.25e-12}} = \num{1.3e-4} \]

The chromate is \textbf{ten times more soluble} despite a $K_{sp}$
eighteen times smaller. A bare comparison of the constants gives the
\emph{wrong} answer here, which is precisely why exam writers love this
pairing.

\textbf{Why the inversion happens.} For a 1:1 salt $s$ scales as
$K_{sp}^{1/2}$; for a 2:1 salt it scales as $K_{sp}^{1/3}$. A cube
root climbs much faster than a square root at these tiny values, so a
2:1 salt squeezes more dissolution out of a smaller constant. You do
not need that sentence on the exam --- you need the habit it protects:
\textbf{different shapes, do the arithmetic.}
\end{workedexample}

\begin{yourturn}[YOUR TURN 4 --- four questions]
\begin{enumerate}[leftmargin=*,itemsep=0.5em]
  \item Rank \ce{AgCl} ($K_{sp} = \num{1.6e-10}$), \ce{AgBr}
        ($\num{5.0e-13}$), \ce{AgI} ($\num{1.5e-16}$) by solubility.
        Justify the shortcut. \workspace[1.8cm]
  \item May you rank \ce{BaSO4} ($\num{1.5e-9}$) against \ce{CaF2}
        ($\num{4.0e-11}$) by comparing the constants? Why or why not?
        \workspace[1.8cm]
  \item Compute both molar solubilities from question 2 and give the
        ranking. \workspace[2.2cm]
  \item A student concludes \ce{Ag2CrO4} is less soluble than \ce{AgCl}
        ``because its $K_{sp}$ is smaller''. What went wrong?
        \workspace[1.8cm]
\end{enumerate}
\selfcheck{(a) \ce{AgCl} $>$ \ce{AgBr} $>$ \ce{AgI} --- same type \quad
(b) no --- different types \quad (c) \ce{CaF2} ($\num{2.2e-4}$) $>$
\ce{BaSO4} ($\num{3.9e-5}$) \quad (d) compared constants across
formula types}
\keyonly{\begin{apcanswer}(a) $\ce{AgCl} > \ce{AgBr} > \ce{AgI}$. All
three are \textbf{1:1 salts}, so $s = \sqrt{K_{sp}}$ for each and the
$K_{sp}$ order is the solubility order. (b) \textbf{No.} \ce{BaSO4} is
1:1 but \ce{CaF2} is 1:2 --- different algebra shapes, so the constants
are not directly comparable. (c) $s(\ce{BaSO4}) = \sqrt{\num{1.5e-9}} =
\num{3.9e-5}$; $s(\ce{CaF2}) = \sqrt[3]{\num{4.0e-11}/4} =
\sqrt[3]{\num{1.0e-11}} = \num{2.2e-4}$. \textbf{\ce{CaF2} is more
soluble} --- despite the smaller $K_{sp}$, the same inversion as the
worked example. (d) They \textbf{compared $K_{sp}$ values across
different formula types}. \ce{Ag2CrO4} is 2:1 and \ce{AgCl} is 1:1;
computing $s$ shows the chromate is ten times \emph{more}
soluble.\end{apcanswer}}
\end{yourturn}

% =====================================================================
\section*{\sffamily Ladder 5 \textbullet\ The common-ion effect}
\zsec{16.1}

Dissolve \ce{AgCl} not in pure water but in salt water --- water that
already contains \ce{Cl-}. Le Ch\^{a}telier says the equilibrium is
pushed back toward the solid, and the arithmetic says by how much.

\begin{center}
\begin{tikzpicture}[font=\sffamily\footnotesize]
  % beaker 1: pure water
  \node[font=\sffamily\small\bfseries] at (1.85,3.3) {in pure water};
  \fill[apcgray!22] (0.6,0.5) rectangle (3.1,2.7);
  \draw[very thick] (0.6,2.95) -- (0.6,0.5) -- (3.1,0.5) -- (3.1,2.95);
  \fill[apcgray!75] (1.3,0.5) rectangle (2.4,0.78);
  \foreach \x/\y in {1.1/1.5, 2.0/2.1, 2.7/1.3, 1.6/1.05} {
    \fill (\x,\y) circle (0.06);}
  \node[align=center] at (1.85,-0.2)
    {$s = \num{1.3e-5}$};

  % beaker 2: 0.10 M NaCl
  \begin{scope}[xshift=4.6cm]
    \node[font=\sffamily\small\bfseries] at (1.85,3.3)
      {in \SI{0.10}{\Molar} \ce{NaCl}};
    \fill[apcgray!22] (0.6,0.5) rectangle (3.1,2.7);
    \draw[very thick] (0.6,2.95) -- (0.6,0.5) -- (3.1,0.5)
      -- (3.1,2.95);
    \fill[apcgray!75] (1.0,0.5) rectangle (2.8,0.85);
    \fill (1.85,1.6) circle (0.06);
    \node[align=center] at (1.85,-0.2)
      {$s = \num{1.6e-9}$};
  \end{scope}

  \node[anchor=west,align=left,font=\sffamily\footnotesize]
    at (9.0,1.6)
    {the added \ce{Cl-} pushes\\
     $\ce{AgCl <=> Ag+ + Cl-}$\\
     back to the \textbf{left}:\\
     8000 times less dissolves};
\end{tikzpicture}
\end{center}

\begin{workedexample}[I do: solubility with a common ion]
\textbf{Find the molar solubility of \ce{AgCl}
($K_{sp} = \num{1.6e-10}$) in \SI{0.10}{\Molar} \ce{NaCl}.}

\textbf{Set up the equilibrium with the chloride head start.} Let $s$
dissolve. Then $[\ce{Ag+}] = s$, but chloride starts at 0.10:
\[ K_{sp} = [\ce{Ag+}][\ce{Cl-}] = s\,(0.10 + s) \]

\textbf{Approximate, Chapter 13 style.} $s$ will be tiny against 0.10
(the $K_{sp}$ is $10^{-10}$), so $0.10 + s \approx 0.10$:
\[ \num{1.6e-10} = s \times 0.10
   \qquad\Longrightarrow\qquad
   s = \mathbf{\SI{1.6e-9}{\Molar}} \]
Check: $\num{1.6e-9}/0.10 \times 100$ is far below 5\% --- the
approximation is comfortably valid.

\textbf{Compare with pure water.} There $s = \num{1.3e-5}$; here
$s = \num{1.6e-9}$ --- the common ion suppressed the solubility by a
factor of about \textbf{8000}. Qualitatively that is Le Ch\^{a}telier
(added product pushes the equilibrium left); quantitatively it falls
straight out of the $K_{sp}$ expression, because $[\ce{Cl-}]$ is pinned
at 0.10 and $[\ce{Ag+}]$ must shrink until the product is back to
$K_{sp}$.

\textbf{The setup difference is the whole skill.} In pure water both
ions are written in terms of $s$. With a common ion, one concentration
has a \emph{head start} that dwarfs $s$. Notice which ion the added
salt supplies, give it the head start, and the rest is Chapter 13.

\textbf{Both directions of the CED skill.} 7.12 also asks it backwards:
given the measured solubility in a solution of known common-ion
concentration, recover $K_{sp} = s \times c$. Same equation, solved for
the other symbol.
\end{workedexample}

\begin{yourturn}[YOUR TURN 5 --- four questions]
\begin{enumerate}[leftmargin=*,itemsep=0.5em]
  \item In which does \ce{BaSO4} dissolve more:\ pure water or
        \SI{0.10}{\Molar} \ce{Na2SO4}? Explain with
        Le Ch\^{a}telier. \workspace[1.8cm]
  \item Find the molar solubility of \ce{BaSO4}
        ($K_{sp} = \num{1.5e-9}$) in \SI{0.10}{\Molar} \ce{Na2SO4}.
        \workspace[2.0cm]
  \item Find the molar solubility of \ce{AgBr}
        ($K_{sp} = \num{5.0e-13}$) in \SI{0.010}{\Molar} \ce{NaBr}.
        \workspace[2.0cm]
  \item Does the common ion change $K_{sp}$? Explain.
        \workspace[1.6cm]
\end{enumerate}
\selfcheck{(a) pure water --- added sulfate pushes left \quad (b)
\SI{1.5e-8}{\Molar} \quad (c) \SI{5.0e-11}{\Molar} \quad (d) no ---
only temperature changes $K_{sp}$}
\keyonly{\begin{apcanswer}(a) \textbf{Pure water.} In \ce{Na2SO4}
solution the added \ce{SO4^2-} is a product of the dissolution
equilibrium, so the system shifts \textbf{left} and less dissolves.
(b) $s \times 0.10 = \num{1.5e-9}$, so $s =
\mathbf{\SI{1.5e-8}{\Molar}}$ (against $\num{3.9e-5}$ in pure water).
(c) $s \times 0.010 = \num{5.0e-13}$, so $s =
\mathbf{\SI{5.0e-11}{\Molar}}$. (d) \textbf{No.} $K_{sp}$ is an
equilibrium constant and changes only with \textbf{temperature}. The
common ion moves the \emph{position} of the equilibrium --- how much
dissolves --- while the product $[\ce{Ag+}][\ce{Cl-}]$ at saturation
stays exactly $K_{sp}$.\end{apcanswer}}
\end{yourturn}

% =====================================================================
\section*{\sffamily Ladder 6 \textbullet\ pH and solubility ---
qualitative by decree}
\zsec{16.1}

Some salts dissolve better in acid. Which ones, and why, is assessed;
\emph{how much better} is explicitly not.

\begin{apcnote}[AP scope: this ladder is deliberately non-numerical]
CED topic 8.11's exclusion: \emph{``Computations of solubility as a
function of pH will not be assessed on the AP Exam.''} What is
assessed is the \textbf{qualitative} call --- identifying \emph{whether}
a salt's solubility is pH-sensitive and \emph{which way} it moves,
argued through Le Ch\^{a}telier. That is exactly what this ladder
practises, and no more.
\end{apcnote}

\begin{center}
\begin{tikzpicture}[font=\sffamily\footnotesize]
  \node[font=\sffamily\small\bfseries] at (5.4,2.9)
    {is the solubility pH-sensitive? look at the anion};
  \foreach \y/\an/\verdict in {%
      2.2/{\ce{OH-}, \ce{F-}, \ce{CO3^2-}, \ce{PO4^3-}, \ce{S^2-}
        --- basic anions}/{\textbf{yes}: acid consumes the anion,
        solubility \textbf{rises}},
      1.45/{\ce{Cl-}, \ce{Br-}, \ce{I-}, \ce{NO3-} --- conjugates of
        strong acids}/{\textbf{no}: the anion ignores \ce{H+}}} {
    \node[anchor=west] at (0.4,\y) {\an};
    \node[anchor=west] at (0.9,\y-0.38) {\verdict};}
  \draw[apcgray] (0.3,0.75) -- (10.8,0.75);
  \node[anchor=west] at (0.4,0.3)
    {test: is the anion the conjugate base of a \emph{weak} acid? then
     \ce{H+} can pull it out of the equilibrium};
\end{tikzpicture}
\end{center}

\begin{workedexample}[I do: why acid dissolves limestone but not
silver chloride]
\textbf{\ce{CaCO3} in acid.} The dissolution equilibrium is
$\ce{CaCO3(s) <=> Ca^2+ + CO3^2-}$. Carbonate is the conjugate base of
the weak acid \ce{HCO3-}, so added \ce{H+} reacts with it:
$\ce{CO3^2- + H+ -> HCO3-}$. That \textbf{removes a product} of the
dissolution equilibrium, and by Le Ch\^{a}telier the system shifts
right --- more limestone dissolves. This is why acid rain eats marble
statues and why geologists drip \ce{HCl} on rocks to identify
carbonates by the fizz.

\textbf{\ce{AgCl} in acid.} Chloride is the conjugate base of the
\emph{strong} acid \ce{HCl} --- which is to say, essentially no base at
all. Added \ce{H+} does not react with it, nothing is removed from the
equilibrium, and the solubility does not change. \ce{AgCl} is as
insoluble in acid as in water.

\textbf{The one-sentence exam answer.} ``The anion is the conjugate
base of a weak acid, so \ce{H+} reacts with it and removes it from
solution; by Le Ch\^{a}telier the dissolution equilibrium shifts toward
the ions, increasing solubility.'' Name the anion, name the removal,
name the shift --- three clauses, full credit.

\textbf{\ce{Mg(OH)2} is the cleanest case of all:} the anion \emph{is}
hydroxide, and \ce{H+ + OH- -> H2O} is the most familiar removal
reaction in the course. Milk of magnesia dissolving in stomach acid is
this ladder in your medicine cabinet.
\end{workedexample}

\begin{yourturn}[YOUR TURN 6 --- four questions]
\begin{enumerate}[leftmargin=*,itemsep=0.5em]
  \item Is \ce{CaF2} more soluble in acid than in water? Explain.
        \workspace[1.8cm]
  \item Is \ce{AgBr} more soluble in acid than in water? Explain.
        \workspace[1.8cm]
  \item Write the reaction by which \ce{H+} increases the solubility of
        \ce{Mg(OH)2}. \workspace[1.6cm]
  \item What single question about the anion decides pH sensitivity?
        \workspace[1.6cm]
\end{enumerate}
\selfcheck{(a) yes --- \ce{F-} is a weak-acid conjugate \quad (b) no
--- \ce{Br-} is a strong-acid conjugate \quad (c)
$\ce{H+ + OH- -> H2O}$ \quad (d) is it the conjugate base of a weak
acid?}
\keyonly{\begin{apcanswer}(a) \textbf{Yes.} \ce{F-} is the conjugate
base of the weak acid \ce{HF}, so \ce{H+} converts it to \ce{HF},
removing a product of the dissolution equilibrium; the system shifts
right and more dissolves. (b) \textbf{No.} \ce{Br-} is the conjugate
base of the strong acid \ce{HBr} --- far too weak a base to react with
\ce{H+} --- so nothing is removed and the solubility is unchanged.
(c) $\ce{Mg(OH)2(s) <=> Mg^2+ + 2OH-}$ followed by
$\mathbf{\ce{H+ + OH- -> H2O}}$ --- the acid consumes the hydroxide and
the equilibrium shifts right. (d) \textbf{Is the anion the conjugate
base of a weak acid?} If yes, \ce{H+} can remove it and solubility
rises at low pH; if it comes from a strong acid, pH has no
effect.\end{apcanswer}}
\end{yourturn}

% =====================================================================
\section*{\sffamily Ladder 7 \textbullet\ $Q$ versus $K_{sp}$ --- will
a precipitate form?}
\zsec{16.2}

Chapter 13's reaction quotient, pointed at dissolution. Compute the ion
product from the concentrations you actually have and compare it with
$K_{sp}$.

\begin{center}
\begin{tikzpicture}[font=\sffamily\footnotesize]
  \draw[-{Stealth[length=3mm]},very thick] (0.6,2.3) -- (10.7,2.3);
  \node[font=\sffamily\bfseries] at (5.6,2.75) {$K_{sp}$};
  \fill (5.6,2.3) circle (0.11);
  \draw[densely dashed,apcgray] (5.6,2.02) -- (5.6,2.25);
  \node[font=\sffamily\bfseries] at (2.5,1.85) {$Q < K_{sp}$};
  \node[align=center] at (2.5,1.1)
    {unsaturated\\\textbf{no precipitate};\\more solid could dissolve};
  \node[font=\sffamily\bfseries] at (5.6,1.85) {$Q = K_{sp}$};
  \node[align=center] at (5.6,1.1) {saturated\\at\\equilibrium};
  \node[font=\sffamily\bfseries] at (8.8,1.85) {$Q > K_{sp}$};
  \node[align=center] at (8.8,1.1)
    {supersaturated\\\textbf{precipitate forms}\\until $Q$ falls to
     $K_{sp}$};
\end{tikzpicture}
\end{center}

\begin{workedexample}[I do: three verdicts on one salt]
\textbf{$K_{sp}(\ce{AgCl}) = \num{1.6e-10}$. For each mixture, decide
what happens.}

\textbf{(a) $[\ce{Ag+}] = \num{1.0e-6}$, $[\ce{Cl-}] =
\num{1.0e-6}$.}
\[ Q = (\num{1.0e-6})(\num{1.0e-6}) = \num{1.0e-12} \]
$Q < K_{sp}$: \textbf{no precipitate}. The solution is unsaturated ---
if solid \ce{AgCl} were present it would keep dissolving.

\textbf{(b) $[\ce{Ag+}] = \num{4.0e-5}$, $[\ce{Cl-}] =
\num{4.0e-5}$.}
\[ Q = \num{1.6e-9} \]
$Q > K_{sp}$ by a factor of ten: \textbf{a precipitate forms}, and it
keeps forming --- consuming ions --- until the product of what remains
in solution has fallen back to $\num{1.6e-10}$.

\textbf{(c) $[\ce{Ag+}] = \num{1.6e-5}$, $[\ce{Cl-}] =
\num{1.0e-5}$.}
\[ Q = \num{1.6e-10} = K_{sp} \]
Exactly saturated: the solution sits at equilibrium, holding as much as
it ever will. Note the two concentrations are \emph{not equal} --- only
their product matters.

\textbf{Two habits from Chapter 13 carry over unchanged.} First, $Q$
uses the same expression as $K$ --- coefficients as exponents, solids
omitted. Second, the system always moves so as to push $Q$ toward
$K_{sp}$: too big, precipitate (removes ions, lowers $Q$); too small,
dissolve (adds ions, raises $Q$). You never memorise which way ---
you read it off the fraction.
\end{workedexample}

\begin{yourturn}[YOUR TURN 7 --- four questions]
$K_{sp}(\ce{BaSO4}) = \num{1.5e-9}$.
\begin{enumerate}[leftmargin=*,itemsep=0.5em]
  \item $[\ce{Ba^2+}] = \num{1.0e-5}$ and $[\ce{SO4^2-}] =
        \num{1.0e-5}$. Precipitate? \workspace[1.8cm]
  \item $[\ce{Ba^2+}] = \num{1.0e-4}$ and $[\ce{SO4^2-}] =
        \num{1.0e-4}$. Precipitate? \workspace[1.8cm]
  \item For \ce{PbI2} ($K_{sp} = \num{1.4e-8}$):
        $[\ce{Pb^2+}] = \num{1.0e-5}$, $[\ce{I-}] = \num{1.0e-3}$.
        Precipitate? (Mind the exponent.) \workspace[2.0cm]
  \item What physically happens, and to which quantity, when
        $Q > K_{sp}$? \workspace[1.8cm]
\end{enumerate}
\selfcheck{(a) $Q = \num{1.0e-10}$, no \quad (b) $Q = \num{1.0e-8}$,
yes \quad (c) $Q = \num{1.0e-11}$, no \quad (d) ions leave solution as
solid until $Q$ falls to $K_{sp}$}
\keyonly{\begin{apcanswer}(a) $Q = (\num{1.0e-5})^{2} =
\num{1.0e-10} < \num{1.5e-9}$: \textbf{no precipitate}. (b)
$Q = (\num{1.0e-4})^{2} = \num{1.0e-8} > \num{1.5e-9}$:
\textbf{precipitate forms}. (c) $Q = [\ce{Pb^2+}][\ce{I-}]^{2} =
(\num{1.0e-5})(\num{1.0e-3})^{2} = \num{1.0e-11} < \num{1.4e-8}$:
\textbf{no precipitate} --- squaring the iodide is the step the
question is testing. (d) Ions pair up and \textbf{leave the solution as
solid} --- precipitation --- which lowers the ion concentrations and
therefore lowers $Q$, continuing until $Q = K_{sp}$ and the system is
at equilibrium.\end{apcanswer}}
\end{yourturn}

% =====================================================================
\section*{\sffamily Ladder 8 \textbullet\ The mixing problem}
\zsec{16.2}

The exam's favourite disguise for Ladder 7: two solutions are poured
together. Mixing \textbf{dilutes both} before any chemistry happens,
and forgetting the dilution is the built-in trap.

\begin{center}
\begin{tikzpicture}[font=\sffamily\footnotesize]
  \node[font=\sffamily\small\bfseries] at (5.6,3.0)
    {the three-step routine};
  \foreach \y/\n/\t in {%
      2.3/{1}/{\textbf{dilute}: new concentration $= M_{1}V_{1} /
        V_{\text{total}}$ for each ion (Ch 11, Ladder 2)},
      1.7/{2}/{\textbf{compute} $Q$ from the diluted concentrations},
      1.1/{3}/{\textbf{compare} with $K_{sp}$ and state the verdict}} {
    \node[anchor=west,font=\sffamily\bfseries] at (0.5,\y) {\n.};
    \node[anchor=west] at (1.05,\y) {\t};}
  \draw[apcgray] (0.4,0.75) -- (10.8,0.75);
  \node[anchor=west,font=\sffamily\itshape] at (0.5,0.3)
    {equal volumes mixed $\Rightarrow$ every concentration simply
     \textbf{halves}};
\end{tikzpicture}
\end{center}

\begin{workedexample}[I do: mix, dilute, decide]
\textbf{\SI{100.}{\milli\litre} of \SI{4.0e-4}{\Molar} \ce{AgNO3} is
mixed with \SI{100.}{\milli\litre} of \SI{4.0e-4}{\Molar} \ce{NaCl}.
Does \ce{AgCl} precipitate? ($K_{sp} = \num{1.6e-10}$)}

\textbf{Step 1 --- dilution.} The total volume is
\SI{200.}{\milli\litre}, so each ion's concentration halves:
\[ [\ce{Ag+}] = [\ce{Cl-}]
   = \num{4.0e-4} \times \frac{100.}{200.}
   = \num{2.0e-4}\;\si{\Molar} \]

\textbf{Step 2 --- compute $Q$.}
\[ Q = (\num{2.0e-4})(\num{2.0e-4}) = \num{4.0e-8} \]

\textbf{Step 3 --- compare.}
\[ \num{4.0e-8} > \num{1.6e-10}
   \qquad\Longrightarrow\qquad
   \textbf{a precipitate forms} \]
$Q$ exceeds $K_{sp}$ by a factor of 250, so most of the silver ends up
as solid \ce{AgCl}.

\textbf{The trap, run deliberately.} Skip the dilution and you compute
$Q = (\num{4.0e-4})^{2} = \num{1.6e-7}$ --- still a ``yes'' here, but
four times too large, and on a question built closer to the boundary
the undiluted $Q$ gives the \emph{wrong verdict}. The dilution step is
where the marks live; distractors are manufactured from exactly this
omission.

\textbf{Sig-fig sanity.} Everything here is two significant figures,
and the comparison spans two orders of magnitude --- so no rounding
choice can flip the answer. Exam items are built with that same
robustness, which is why a clean setup matters more than calculator
precision.
\end{workedexample}

\begin{yourturn}[YOUR TURN 8 --- four questions]
\begin{enumerate}[leftmargin=*,itemsep=0.5em]
  \item Equal volumes of \SI{2.0e-3}{\Molar} \ce{Pb(NO3)2} and
        \SI{2.0e-3}{\Molar} \ce{KI} are mixed. Give the diluted
        $[\ce{Pb^2+}]$ and $[\ce{I-}]$. \workspace[1.8cm]
  \item Compute $Q$ for \ce{PbI2} from those diluted values.
        \workspace[2.0cm]
  \item $K_{sp}(\ce{PbI2}) = \num{1.4e-8}$. Does a precipitate form?
        \workspace[1.6cm]
  \item Why must the dilution come \emph{before} the $Q$ calculation?
        \workspace[1.8cm]
\end{enumerate}
\selfcheck{(a) \num{1.0e-3} each \quad (b) $Q = \num{1.0e-9}$ \quad
(c) no --- $Q < K_{sp}$ \quad (d) mixing changes the concentrations
$Q$ must be computed from}
\keyonly{\begin{apcanswer}(a) Equal volumes halve both:
$[\ce{Pb^2+}] = [\ce{I-}] = \mathbf{\num{1.0e-3}}\;\si{\Molar}$.
(b) $Q = [\ce{Pb^2+}][\ce{I-}]^{2} = (\num{1.0e-3})(\num{1.0e-3})^{2}
= \mathbf{\num{1.0e-9}}$. (c) $\num{1.0e-9} < \num{1.4e-8}$:
\textbf{no precipitate} --- despite both solutions sounding
``concentrated enough'', the cube of a millimolar number is small.
Without the dilution step the verdict would flip: $(\num{2.0e-3})^{3}
= \num{8.0e-9}$\ldots\ still under, but $Q$ computed wrongly by a
factor of 8. (d) Because $Q$ is defined by the concentrations
\textbf{in the actual mixture}, and pouring the solutions together
changed every concentration. Computing $Q$ from the bottle
concentrations describes a solution that no longer
exists.\end{apcanswer}}
\end{yourturn}

% =====================================================================
\section*{\sffamily Ladder 9 \textbullet\ Reading and representing
saturation}
\zsec{16.1--16.2}

CED 7.8's ``representations'' skill, applied to solubility: particulate
diagrams, and the vocabulary of saturation, said precisely.

\begin{center}
\begin{tikzpicture}[font=\sffamily\footnotesize]
  \foreach \i/\nm/\desc in {%
      0/{unsaturated}/{$Q < K_{sp}$\\could dissolve more\\no solid
        persists},
      4.2/{saturated}/{$Q = K_{sp}$\\dynamic equilibrium\\solid may sit
        at the bottom},
      8.4/{supersaturated}/{$Q > K_{sp}$\\unstable\\will precipitate}} {
    \begin{scope}[xshift=\i cm]
      \node[font=\sffamily\small\bfseries] at (1.5,3.4) {\nm};
      \fill[apcgray!22] (0.55,0.5) rectangle (2.45,2.6);
      \draw[very thick] (0.55,2.85) -- (0.55,0.5) -- (2.45,0.5)
        -- (2.45,2.85);
      \node[align=center] at (1.5,-0.45) {\desc};
    \end{scope}}
  % unsaturated: few ions, no solid
  \foreach \p in {(0.95,1.2),(1.6,1.9),(2.1,1.0),(1.2,2.2)} {
    \fill \p circle (0.06);}
  % saturated: ions + solid
  \begin{scope}[xshift=4.2cm]
    \foreach \p in {(0.9,1.3),(1.5,2.0),(2.1,1.1),(1.2,1.7),
                    (1.9,2.25),(1.55,1.05)} {
      \fill \p circle (0.06);}
    \fill[apcgray!75] (0.95,0.5) rectangle (2.05,0.78);
  \end{scope}
  % supersaturated: crowded, no solid yet
  \begin{scope}[xshift=8.4cm]
    \foreach \p in {(0.85,1.1),(1.3,1.5),(1.8,1.05),(2.15,1.6),
                    (1.0,1.95),(1.55,2.25),(2.05,2.1),(1.35,0.85),
                    (0.8,2.35),(1.85,1.85)} {
      \fill \p circle (0.06);}
  \end{scope}
\end{tikzpicture}
\end{center}

\begin{workedexample}[I do: the sentences that earn credit]
\textbf{FRQ prompts in this unit ask you to \emph{explain}, and partial
sentences lose partial credit. Four statements, said exactly.}

\textbf{What ``saturated'' means.} ``The solution is in dynamic
equilibrium with the solid: dissolution and precipitation continue at
equal rates, so the ion concentrations remain constant and
$[\ce{Ag+}][\ce{Cl-}] = K_{sp}$.'' The word \emph{dynamic} --- and the
equal-rates clause backing it --- is the same mark as Chapter 13,
Ladder 1.

\textbf{Why adding more solid changes nothing.} ``The solid does not
appear in the equilibrium expression; as long as some solid is present,
adding more changes neither the ion concentrations nor the position of
equilibrium.'' (Chapter 13, Ladder 4 --- now in its natural habitat.)

\textbf{What a particulate diagram must show for \emph{saturated}.}
Free ions in the ratio of the formula --- two \ce{Ag+} drawn for every
\ce{CrO4^2-} if the salt is \ce{Ag2CrO4} --- with excess solid drawn as
a lattice at the bottom. Drawing intact ``\ce{AgCl} molecules''
floating in solution is the standard wrong answer: a dissolved ionic
compound is \emph{ions}, full stop.

\textbf{What happens on evaporation.} ``Removing water raises the ion
concentrations, so $Q$ rises above $K_{sp}$ and solid precipitates
until $Q$ returns to $K_{sp}$.'' Evaporation questions are Ladder 7
wearing a lab coat: the comparison, not the water, is the chemistry.
\end{workedexample}

\begin{yourturn}[YOUR TURN 9 --- four questions]
\begin{enumerate}[leftmargin=*,itemsep=0.5em]
  \item A saturated \ce{AgCl} solution sits over excess solid. Are
        dissolution and precipitation still occurring? Explain.
        \workspace[1.8cm]
  \item In a particulate diagram of saturated \ce{CaF2} solution, what
        ratio of free ions must appear? \workspace[1.6cm]
  \item Half the water in a saturated \ce{BaSO4} solution (no excess
        solid) evaporates. What happens, in $Q$-versus-$K_{sp}$
        language? \workspace[2.0cm]
  \item Why is drawing dissolved ``\ce{NaCl} units'' in a particulate
        diagram wrong? \workspace[1.6cm]
\end{enumerate}
\selfcheck{(a) yes --- at equal rates \quad (b) two \ce{F-} per
\ce{Ca^2+} \quad (c) concentrations double, $Q > K_{sp}$, precipitate
forms \quad (d) dissolved ionic compounds exist as separate ions}
\keyonly{\begin{apcanswer}(a) \textbf{Yes.} Equilibrium is dynamic:
ions leave the crystal and redeposit on it continuously, at
\textbf{equal rates}, which is why nothing appears to change. (b)
\textbf{Two fluorides for every calcium} --- the diagram must show the
formula's stoichiometry among the free ions. (c) Halving the water
\textbf{doubles both concentrations}, so $Q = (2[\ce{Ba^2+}])
(2[\ce{SO4^2-}]) = 4K_{sp} > K_{sp}$ --- solid precipitates until the
remaining ion product falls back to $K_{sp}$. (d) Because a dissolved
ionic compound exists as \textbf{separate, independent hydrated ions},
not as neutral pairs. Particulate diagrams are graded on exactly this
point.\end{apcanswer}}
\end{yourturn}

% =====================================================================
\section*{\sffamily Mixed set \textbullet\ ten questions, no labels}

\begin{apcnote}[Why this section exists]
Every question so far told you which ladder it belonged to. These ten
do not --- deciding \emph{which} skill a solubility question wants is
most of the difficulty on the real exam.
\end{apcnote}

\begin{yourturn}[MIXED SET --- first five]
\begin{enumerate}[leftmargin=*,itemsep=0.5em,label=(\alph*)]
  \item Write the $K_{sp}$ expression for \ce{Fe(OH)3}.
        \workspace[1.8cm]
  \item A 1:1 salt has molar solubility \SI{2.0e-4}{\Molar}. Find
        $K_{sp}$. \workspace[1.8cm]
  \item $K_{sp}(\ce{MX2}) = \num{4.0e-9}$. Find the molar solubility.
        \workspace[2.0cm]
  \item Salt A (1:1) has $K_{sp} = \num{1.0e-10}$; salt B (1:2) has
        $K_{sp} = \num{1.0e-11}$. May you conclude A is more soluble?
        \workspace[1.8cm]
  \item Is \ce{SrCO3} more soluble at pH 4 than at pH 7? One sentence.
        \workspace[1.8cm]
\end{enumerate}
\selfcheck{(a) $[\ce{Fe^3+}][\ce{OH-}]^{3}$ \quad (b) \num{4.0e-8}
\quad (c) \SI{1.0e-3}{\Molar} \quad (d) no --- different formula
types \quad (e) yes --- carbonate is a weak-acid conjugate}
\keyonly{\begin{apcanswer}\textbf{(a)} $K_{sp} =
[\ce{Fe^3+}][\ce{OH-}]^{3}$ --- coefficient 3 becomes the exponent.
\textbf{(b)} $K_{sp} = s^{2} = (\num{2.0e-4})^{2} =
\mathbf{\num{4.0e-8}}$. \textbf{(c)} $4s^{3} = \num{4.0e-9}$, so
$s^{3} = \num{1.0e-9}$ and $s = \mathbf{\SI{1.0e-3}{\Molar}}$.
\textbf{(d)} \textbf{No} --- the formulas have different shapes
($s^{2}$ versus $4s^{3}$), so the constants are not directly
comparable; compute both solubilities. \textbf{(e)} \textbf{Yes}:
\ce{CO3^2-} is the conjugate base of a weak acid, so at pH 4 the
\ce{H+} removes it from solution and the dissolution equilibrium
shifts right.\end{apcanswer}}
\end{yourturn}

\begin{yourturn}[MIXED SET --- last five]
\begin{enumerate}[leftmargin=*,itemsep=0.5em,label=(\alph*)]
  \item Find the molar solubility of \ce{AgCl}
        ($K_{sp} = \num{1.6e-10}$) in \SI{0.20}{\Molar} \ce{KCl}.
        \workspace[2.0cm]
  \item $[\ce{Ca^2+}] = \num{2.0e-4}$ and $[\ce{F-}] = \num{2.0e-4}$
        in a mixture; $K_{sp}(\ce{CaF2}) = \num{4.0e-11}$. Precipitate?
        \workspace[2.0cm]
  \item Equal volumes of \SI{2.0e-4}{\Molar} \ce{AgNO3} and
        \SI{2.0e-4}{\Molar} \ce{NaBr} are mixed;
        $K_{sp}(\ce{AgBr}) = \num{5.0e-13}$. Precipitate?
        \workspace[2.2cm]
  \item Adding solid \ce{NaCl} to a saturated \ce{AgCl} solution makes
        more \ce{AgCl} precipitate. Explain in one sentence.
        \workspace[1.8cm]
  \item A saturated solution of \ce{Ag2CrO4} has $[\ce{CrO4^2-}] =
        \num{1.3e-4}$. What is $[\ce{Ag+}]$, and what $K_{sp}$ does
        this give? \workspace[2.2cm]
\end{enumerate}
\selfcheck{(a) \SI{8.0e-10}{\Molar} \quad (b) $Q = \num{8.0e-12}$, no
\quad (c) $Q = \num{1.0e-8}$, yes \quad (d) common ion pushes the
equilibrium left \quad (e) \num{2.6e-4}; $K_{sp} \approx \num{8.8e-12}$}
\keyonly{\begin{apcanswer}\textbf{(a)} The chloride head start is 0.20,
so $s \times 0.20 = \num{1.6e-10}$ and $s =
\mathbf{\SI{8.0e-10}{\Molar}}$. \textbf{(b)} $Q =
(\num{2.0e-4})(\num{2.0e-4})^{2} = \num{8.0e-12} < \num{4.0e-11}$:
\textbf{no precipitate}. \textbf{(c)} Halve both: \num{1.0e-4} each;
$Q = \num{1.0e-8} > \num{5.0e-13}$: \textbf{precipitate forms} --- by
four orders of magnitude. \textbf{(d)} The added \ce{Cl-} is a common
ion, so the product $[\ce{Ag+}][\ce{Cl-}]$ exceeds $K_{sp}$ and the
equilibrium shifts left, precipitating \ce{AgCl} until the product
returns to $K_{sp}$. \textbf{(e)} The formula gives two silvers per
chromate: $[\ce{Ag+}] = 2 \times \num{1.3e-4} =
\mathbf{\num{2.6e-4}}$; then $K_{sp} = (\num{2.6e-4})^{2}
(\num{1.3e-4}) = \mathbf{\num{8.8e-12}}$ --- close to the tabulated
$\num{9.0e-12}$, the gap being rounding in
$s$.\end{apcanswer}}
\end{yourturn}

% =====================================================================
\section*{\sffamily Mastery tracker}

Tick a row only if \textbf{all four} YOUR TURN questions were right on
the first attempt. Every row is assessed --- the off-CED parts of this
chapter (selective precipitation, qualitative analysis, complex ions)
were cut rather than included, and Ladder 6 is deliberately qualitative
because the CED excludes the computation.

\renewcommand{\arraystretch}{1.3}
\noindent\begin{tabular}{@{}c p{5.9cm} c p{4.9cm}@{}}
\toprule
\textbf{First try?} & \textbf{Skill} & \textbf{Ladder} &
  \textbf{If not, re-read\ldots}\\
\midrule
$\square$ & writing $K_{sp}$ expressions & 1 & no denominator \\
$\square$ & solubility $\to$ $K_{sp}$ & 2 & doubled AND squared \\
$\square$ & $K_{sp}$ $\to$ solubility & 3 & $s^{2}$ or $4s^{3}$ \\
$\square$ & comparing solubilities & 4 & same type, or compute \\
$\square$ & the common-ion effect & 5 & give one ion a head start \\
$\square$ & pH and solubility & 6 & weak-acid conjugate anions \\
$\square$ & $Q$ versus $K_{sp}$ & 7 & the verdict is a comparison \\
$\square$ & the mixing problem & 8 & dilute first, always \\
$\square$ & representing saturation & 9 & dynamic, ions, ratio \\
\midrule
$\square$ & \textbf{mixed set} (8 of 10) & --- & whichever it exposed \\
\bottomrule
\end{tabular}

\begin{apcnote}[Scoring yourself honestly]
9/9 and the mixed set: solubility equilibria are banked --- and note
how little was genuinely new. Ladders 1 and 9 were Chapter 13 restated;
5 was the small-$x$ approximation; 7 and 8 were $Q$ versus $K$. This
chapter is a reunion, not a new subject.

6--8: if the misses are Ladders 2--4, the issue is the $s$/$2s$/$4s^{3}$
algebra --- redo those three in sequence, writing the ion
concentrations explicitly every time. If the misses are 7--8, go back
to Chapter 13 Ladder 6, because $Q$ reasoning is upstream of both.

5 or fewer: re-read Chapter 13 Ladders 4 and 6 before touching this
chapter again --- every difficulty here is one of those two in
disguise.
\end{apcnote}

\end{document}
